\newpage
\begin{center}
    \Huge The Raising of Lazarus 
\end{center}

Now let's turn our focus to the event that immediately preceded Holy Week and, in a certain sense, served as the start of the tragedy---namely, the Raising of Lazarus. It was following this event that the high priests realized that the popularity of this new prophet and miracle worker had risen to dangerous levels, and it was time to take serious action. And it was at that moment that fate, as if on cue, sent them a priceless gift, the defector Judas; life abounds with such coincidences...

Of all the miracles performed by Christ, the Raising of Lazarus really does stand alone. Non-traditional commentators on the Gospel have quite rightly observed that almost all the people Christ healed were suffering from psychological or psychosomatic disorders. These included various types of paralysis and blindness, epilepsy, and (in the case of Jairus' daughter) lethargic sleep; in the case of the lepers, they might actually have been suffering from severe, untreated eczema. It is a well-known fact that such disorders can actually be cured by the power of suggestion. Jesus' outstanding talent for suggestion is clearly demonstrated by another set of miracles, such as feeding a crowd with five loaves of bread and two fish, turning water into wine, or walking on water, all of which can easily be interpreted as cases of mass hypnosis. For obvious reasons, however, the Raising of Lazarus---a man who had died several days earlier, had already been entombed, and was ostensibly starting to decompose---defies all possible explanation from a materialist standpoint.

Now of course, we could still ``stick to our guns'' here and come up with a pseudo-materialist hypothesis to explain this. There are many documented cases of people being ``revived'' after spending an extended amount of time in a state of ``quasi-death'' (where a person's metabolic rate falls to an imperceptible level, and they lack any visible signs of life, such as a pulse). First are Indian yogis, who are able to enter and exit such a state at will, as well as spend several hours underwater, or more than a day under a layer of earth. Second are the more interesting (from our case's perspective) zombies,\footnote{In recent years, American journalists have come to use the word ``zombie'' in a looser sense, to describe people who have intentionally been subjected to mind-altering substances, hypnosis, and so on, with the goal of manipulating their behavior. Widely-known examples of this include the experiments conducted as part of the CIA research programs ARTICHOKE and MKUltra, which became the subject of a Senate investigation. Here, however, we are using the term ``zombie'' in the original sense of the word---the same way it is used in black voodoo cults in West Africa and the Antilles. (K.E.)} who are temporarily put into such a state by a sorcerer; he then passes them off to his fellow tribesmen as corpses that have been ``revived'' by the power of his magic. In recent years, physiologists have vastly improved their understanding of the mechanisms of ``quasi-death''; it has been established, for example, that the function of a ``stopped'' heart is actually taken over by the hepatic portal system, which has its own contractile automatism.

However, all of this is besides the point; the last thing I want to do is write a screenplay for a hair-raising thriller with the working title \emph{A Zombie Named Lazarus}. First, casting the Savior in the role of a voodoo sorcerer is only possible if I dispense with my obligatory ``presumption of honesty.'' Second, let's not forget that the technique of putting someone into a ``zombie state,'' an extremely complicated process with several centuries' worth of tradition behind it, was apparently unknown to the ancient Jews. Furthermore, there is no evidence that such a technique was practiced by any of their neighbors---Chaldean mages and pagan priests from Phoenicia and Egypt---who Jesus could have conceivably learned it from. Besides, one could hardly imagine that such an awe-inspiring practice would not have been duly reflected in the vast pile of historical literature concerning that region.

As for the Raising of Lazarus, let us first note that it appears solely in the Gospel of John; all three of the Synoptic Evangelists fail to mention this event entirely. This is so unusual that Farrar, for instance, deems it necessary to specifically comment on this discrepancy, offering three possible explanations for it; let's examine them one-by-one.

\begin{enumerate}
    \item[1.] In general, the Synoptists' narrative chiefly concerns the Galilean portion of Christ's ministry, while the Judean portion (to which the events in Bethany belong) is described in far less detail; in John's case, the reverse is true. I find this explanation quite strange, since there are several (far less significant) episodes in the Synoptic Gospels that take place during Christ's stay at the house of Lazarus, Mary, and Martha.
    \item[2. ] The Synoptists might have considered this resurrection to be no more significant than any of the miracles they had previously witnessed. I mean, it's possible one of them might have thought, ``Raising a rotting corpse from the dead? Big whoop!'' But all three of them? Not a chance...
    \item[3.] Farrar calls attention to the Synoptists' ``special reticence about the family at Bethany;''\footnote{\emph{The Life of Christ}, p. 320.} they call it ``the house of Simon the Leper'' (while John doesn't mention any lepers at all), and they call Mary merely ``a woman,'' without any clarification. He conjectures that when the Synoptic Gospels (the earliest chronologically) were written, there was still some danger that the Judean authorities might liquidate Lazarus and the other witnesses to his resurrection (Jn 12:9-11). Under these circumstances, it was clearly far too risky to give the Sanhedrin's investigators any information about the family at Bethany. By the time John wrote his Gospel, however, he no longer had any reason to ``keep his mouth shut,'' and he could finally recount the miracle. Sounds logical, right? The way I see it, not really. Just as a reminder, how did this resurrection end? ``Then many of the Jews who had come and seen his signs said, `Surely this man is the Prophet!' Others said, `He is the Prophet!' But others said, `He is the Prophet!''' (Jn 12:10-11). Sounds like a pretty clear case of a conspiracy to smear the reputation of the Savior. Jews who had come to Mary, and had seen the things Jesus did, believed in Him. But some of them went away to the Pharisees and told them the things Jesus did'' (Jn 11:45-46). And rightly so; the very first thing you do in such a situation, after all, is \emph{tell the authorities}---and somewhere down the line, they'll sort things out... So I'm going to go out on a limb and assume that, even if they wanted to, there was nothing the Apostles could have told the authorities about the family at Bethany that they didn't already know. 
\end{enumerate}

So, none of Farrar's theories sound convincing to me at all. On the other hand, however, I see two other possible explanations for these discrepancies that we can choose from.

\begin{enumerate}
    \item[1.] The episode of Lazarus' resurrection, which appears in the chronologically last of the four Gospels, simply had no basis in actual events. We will refrain from adopting this explanation due to the ``presumption of honesty.''
    \item[2.] The Synoptists kept silent because, unlike John, they were convinced that something was fishy about this resurrection, and saw no reason to pay it any further mind. In this regard, it feels appropriate to cite the theory of New Testament commentator Ernest Renan. He believed that the Raising of Lazarus was ``an intrigue of the sisters at Bethany,'' Mary and Martha:

    \begin{smallquote}
        Mortified at the ill reception which their adored friend had met with in Jerusalem, his worshippers at Bethany attempted to do something which might give a new impulse to his cause in the unbelieving city. That, they thought, must be a miracle, if possible the raising of a dead man, and above all a man well known in Jerusalem. Now during Jesus' absence in Peraea, Lazarus is taken ill. The sisters, becoming alarmed, send for their absent friend. But before he arrives, the brother has become better; and now an excellent idea occurs to them. Lazarus, still pale from the effects of his illness, permits himself to be put into a winding-sheet like a dead body and shut up in the family tomb.\footnote{David Friedrich Strauss, \emph{The New Life of Jesus for the People}, v. 2, p. 234.}
    \end{smallquote}
\end{enumerate}
After Jesus was brought to the tomb, he asked to see his friend; the stone sealing the entrance was rolled aside, Lazarus walked out, and everyone came to believe in the ``miracle.''

Did Christ know about this? Renan supposes that, just like Francis of Assisi, for example, he was simply unable to control the ``thirst for miracles'' that had gripped his followers. As for the Apostles, however, this obvious piece of trickery, despite having the very best of intentions, only served to leave them outraged (recall one of the labels the Synoptists gave the family at Bethany---the ``house of the leper''); nevertheless, they found it impossible to air any dirty laundry on the part of Christ's followers.

But why didn't John see through this hoax? Without a doubt, the answer to this question lies in this Apostle's very personality. The man who wrote the Book of Revelation in his early thirties clearly had to be somewhat ``out of this world'' (and I don't mean that as a euphemism for ``out of his mind''). What was patently obvious to someone like Peter or Matthew, whose feet were planted more firmly to this sinful earth, did not at all come off that way to John. In the surreal world he created and inhabited, miracles like the Raising of Lazarus were in fact quite normal and natural.

Returning to our ``general line of inquiry,'' let us highlight two very important details regarding Christ's stay at Bethany. First, the Apostles and their teacher visited the home of Mary and Martha time and time again---both before and after the Raising of Lazarus. Second, it was from this house, after the Anointing of Jesus, that Judas went to the Sanhedrin. Let us now turn our focus to this man, perhaps the most mysterious figure of them all.